\subsection{REAL FORMS ASSOCIATED TO A CHEVALLEY SYSTEM}

In this number, we consider a split semi-simple Lie algebra \((\mathfrak{a},\mathfrak{h})\) over the field \(\mathbf{C}\) (Chap. VIII, §2, no. 1), with root system \(\mathrm{R}(\mathfrak{a},\mathfrak{h}) = \mathrm{R}\) , and a Chevalley system \((X_{\alpha})_{\alpha \in \mathbb{R}}\) of \((\mathfrak{a},\mathfrak{h})\) (Chap. VIII, §2, no. 4, Def. 3).

Recall (loc. cit.) that the linear map \(\theta: a \to a\) that coincides with \(-Id_{h}\) on h and maps \(X_{\alpha}\) to \(X_{-\alpha}\) for all \(\alpha \in R\) is an automorphism of a. Moreover (loc. cit., Prop. 7), if \(\alpha, \beta, \alpha + \beta\) are roots, then

\[
[ X _ {\alpha}, X _ {\beta} ] = \mathrm{N} _ {\alpha , \beta} X _ {\alpha + \beta}\tag{3}
\]

with \(\mathrm{N}_{\alpha ,\beta}\in \mathbf{R}^*\) and

\[
\mathrm{N} _ {- \alpha , - \beta} = \mathrm{N} _ {\alpha , \beta}.\tag{4}
\]

Denote by \(\mathfrak{h}_0\) the real vector subspace of \(\mathfrak{h}\) consisting of the \(H\in \mathfrak{h}\) such that \(\alpha (H)\in \mathbf{R}\) for all \(\alpha \in \mathbb{R}\) . Then \(\mathfrak{h}_0\) is an \(\mathbf{R}\) -structure on the complex vector space \(\mathfrak{h}\) , we have \([X_{\alpha},X_{-\alpha}]\in \mathfrak{h}_0\) for all \(\alpha \in \mathbb{R}\) , and the restriction of the Killing form B of \(\mathfrak{a}\) to \(\mathfrak{h}_0\) is separating positive (Chap. VIII, §2, no. 2, Remark 2). Moreover,

\[
\mathrm{B} (H, X _ {\alpha}) = 0, \quad \mathrm{B} (X _ {\alpha}, X _ {\beta}) = 0 \text {if} \alpha + \beta \neq 0, \quad \mathrm{B} (X _ {\alpha}, X _ {- \alpha}) <   0\tag{5}
\]

(Chap. VIII, §2, no. 2, Prop. 1 and no. 4, Lemma 3).

PROPOSITION 2. a) The real vector subspace \(\mathfrak{a}_0 = \mathfrak{h}_0 + \sum_{\alpha \in \mathrm{R}} \mathbf{R}X_\alpha\) of \(\mathfrak{a}\) is a real form of \(\mathfrak{a}\) , of which \(\mathfrak{h}_0\) is a Cartan subalgebra. The pair \((\mathfrak{a}_0, \mathfrak{h}_0)\) is a split semi-simple real Lie algebra, of which \((X_\alpha)\) is a Chevalley system.

\begin{enumerate}
\def\labelenumi{\alph{enumi})}
\setcounter{enumi}{1}
\tightlist
\item
  Let \(\sigma\) be the conjugation of \(\mathfrak{a}\) relative to \(\mathfrak{a}_0\) . Then \(\sigma \circ \theta = \theta \circ \sigma\) . The set of fixed points of \(\sigma \circ \theta\) is a compact real form \(\mathfrak{a}_u\) of \(\mathfrak{a}\) , of which \(i\mathfrak{h}_0\) is a Cartan subalgebra.
\end{enumerate}

Part a) follows immediately from the preceding. We prove b). Since \(\sigma\circ\theta\) and \(\theta\circ\sigma\) are two semi-linear maps from a to a that coincide on \(a_{0}\) , they coincide. Now \(\sigma\circ\theta\) satisfies conditions (2) of no. 1, hence is the conjugation of a relative to the real form \(a_{u}\) consisting of the \(x\in a\) such that \(\sigma\circ\theta(x)=x\) (Prop. 1). For all \(\alpha\in R\) put

\[
u _ {\alpha} = X _ {\alpha} + X _ {- \alpha}, v _ {\alpha} = i (X _ {\alpha} - X _ {- \alpha}).\tag{6}
\]

Then the R-vector space \(a_{u}\) is generated by \(ih_{0}\) , the \(u_{\alpha}\) and the \(v_{\alpha}\) . More precisely, if we choose a chamber C of R, then

\[
\mathfrak {a} _ {u} = i \mathfrak {h} _ {0} \oplus \bigoplus_ {\alpha \in \mathrm{R} _ {+} (\mathrm{C})} \left(\mathbf {R} u _ {\alpha} + \mathbf {R} v _ {\alpha}\right).\tag{7}
\]

It is clear that \(i\mathfrak{h}_0\) is a Cartan subalgebra of \(\mathfrak{a}_u\) , and it remains to prove that the restriction of \(B\) to \(\mathfrak{a}_u\) is negative. Now \(i\mathfrak{h}_0\) and the different subspaces of the form \(\mathbf{R}u_{\alpha}\oplus \mathbf{R}v_{\alpha}\) are orthogonal with respect to B, by (5); the restriction of B to \(i\mathfrak{h}_0\) is negative and

\[
\mathrm{B} (u _ {\alpha}, u _ {\alpha}) = \mathrm{B} (v _ {\alpha}, v _ {\alpha}) = 2 \mathrm{B} (X _ {\alpha}, X _ {- \alpha}) <   0, \quad \mathrm{B} (u _ {\alpha}, v _ {\alpha}) = 0,\tag{8}
\]

hence the conclusion.

Remark. With the preceding notations, we have the following formulas:

\[
[ h, u _ {\alpha} ] = - i \alpha (h) v _ {\alpha}, [ h, v _ {\alpha} ] = i \alpha (h) u _ {\alpha}, [ u _ {\alpha}, v _ {\alpha} ] = 2 i H _ {\alpha}, (h \in \mathfrak {h})\tag{9}
\]

\[
[ u _ {\alpha}, u _ {\beta} ] = \mathrm{N} _ {\alpha , \beta} u _ {\alpha + \beta} + \mathrm{N} _ {\alpha , - \beta} u _ {\alpha - \beta}, \quad \alpha \neq \pm \beta ,\tag{10}
\]

\[
[ v _ {\alpha}, v _ {\beta} ] = - \mathrm{N} _ {\alpha , \beta} u _ {\alpha + \beta} + \mathrm{N} _ {\alpha , - \beta} u _ {\alpha - \beta}, \quad \alpha \neq \pm \beta ,\tag{11}
\]

\[
[ u _ {\alpha}, v _ {\beta} ] = \mathrm{N} _ {\alpha , \beta} v _ {\alpha + \beta} - \mathrm{N} _ {\alpha , - \beta} v _ {\alpha - \beta}, \quad \alpha \neq \pm \beta ,\tag{12}
\]

(in the last three formulas, it is understood, as usual, that \(\mathrm{N}_{\gamma, \delta} = 0\) if \(\gamma + \delta\) is not a root).

Note that \(\sum \mathbf{R}u_{\alpha}\) is a real subalgebra of \(\mathfrak{a}\) , namely \(\mathfrak{a}_0 \cap \mathfrak{a}_u\) .

Let \(\mathrm{Q}(\mathrm{R})\) be the group of radical weights of R (Chap. VI, §1, no. 9). Recall that to any homomorphism \(\gamma : \mathrm{Q}(\mathrm{R}) \to \mathbf{C}^*\) is associated an elementary automorphism \(f(\gamma)\) of \(\mathfrak{a}\) such that \(f(\gamma)(h) = h\) for all \(h \in \mathfrak{h}\) and \(f(\gamma)X_{\alpha} = \gamma(\alpha)X_{\alpha}\) (Chap. VIII, §5, no. 2).

PROPOSITION 3. Let \(\mathfrak{g}\) be a compact real form of \(\mathfrak{a}\) such that \(\mathfrak{g} \cap \mathfrak{h} = i\mathfrak{h}_0\) . There exists a homomorphism \(\gamma: \mathrm{Q}(\mathrm{R}) \to \mathbf{R}_{+}^{*}\) such that \(\mathfrak{g} = f(\gamma)(\mathfrak{a}_u)\) .

Let \(\tau\) be the conjugation of \(\mathfrak{a}\) relative to \(\mathfrak{g}\) . By hypothesis \(\tau(x) = x\) for \(x \in i\mathfrak{h}_0\) , so \(\tau(x) = -x\) for \(x \in \mathfrak{h}_0\) . Thus, for all \(\alpha \in \mathrm{R}\) and all \(h \in \mathfrak{h}_0\) ,

\[
[ h, \tau (X _ {\alpha}) ] = [ - \tau (h), \tau (X _ {\alpha}) ] = - \tau ([ h, X _ {\alpha} ]) = - \tau (\alpha (h) X _ {\alpha});
\]

it follows that \([h, \tau(X_{\alpha})] = -\alpha(h)\tau(X_{\alpha})\) for all \(h \in \mathfrak{h}_0\) , hence also for all \(h \in \mathfrak{h}\) . Hence there exists \(c_{\alpha} \in \mathbf{C}^*\) such that \(\tau(X_{\alpha}) = c_{\alpha}X_{-\alpha}\) . Since \([X_{\alpha}, X_{-\alpha}] \in \mathfrak{h}_0\) , we have \([\tau(X_{\alpha}), \tau(X_{-\alpha})] = -[X_{\alpha}, X_{-\alpha}]\) , so \(c_{\alpha}c_{-\alpha} = 1\) ; similarly, formulas (3) and (4) give \(c_{\alpha + \beta} = c_{\alpha}c_{\beta}\) if \(\alpha, \beta, \alpha + \beta\) are roots. By Chap. VI, §1, no. 6, Cor. 2 of Prop. 19, there exists a homomorphism \(\delta: \mathrm{Q}(\mathrm{R}) \to \mathbf{C}^*\) such that \(\delta(\alpha) = c_{\alpha}\) for all \(\alpha \in \mathrm{R}\) .

We now show that each \(c_{\alpha}\) is strictly positive. Indeed, \(c_{\alpha}\mathrm{B}(X_{\alpha},X_{-\alpha})=\mathrm{B}(X_{\alpha},\tau(X_{\alpha}))\) , and since \(\mathrm{B}(X_{\alpha},X_{-\alpha})\) is negative, it suffices to show that \(\mathrm{B}(z,\tau(z))<0\) for every non-zero element z of a; but every element of a can be written as \(x+iy\) , with x and y in g, and

\[
\mathrm{B} (x + i y, \tau (x + i y)) = \mathrm{B} (x + i y, x - i y) = \mathrm{B} (x, x) + \mathrm{B} (y, y),
\]

hence the stated assertion, the restriction of B to g being negative and separating by hypothesis.

It follows that the homomorphism \(\delta\) takes values in \(\mathbf{R}_+^*\) ; hence there exists a homomorphism \(\gamma : \mathrm{Q}(\mathrm{R}) \to \mathbf{R}_+^*\) such that \(\delta = \gamma^{-2}\) . Then \(f(\gamma)^{-1}(\mathfrak{g})\) is a real form of \(\mathfrak{a}\) ; the corresponding conjugation is \(\tau' = f(\gamma)^{-1} \circ \tau \circ f(\gamma)\) . For all \(\alpha \in \mathbb{R}\) , we have

\[
\tau^ {\prime} (X _ {\alpha}) = f (\gamma) ^ {- 1} (\tau (c _ {\alpha} ^ {- 1 / 2} X _ {\alpha})) = f (\gamma) ^ {- 1} (c _ {\alpha} ^ {1 / 2} X _ {- \alpha}) = X _ {- \alpha},
\]

and \(\tau'(h) = \tau(h) = h\) for \(h \in i\mathfrak{h}_0\) ; it follows that \(\tau'\) is the conjugation with respect to \(\mathfrak{a}_u\) , and hence that \(f(\gamma)^{-1}(\mathfrak{g}) = \mathfrak{a}_u\) .

